% !TeX encoding = UTF-8
% 吉田第六版：中译本第 4 章；此文件当前仅为目录与正文占位。
\SetChapterInfo{subtitle={吉田耕作第六版主线 · 正文待撰写}}
\chapter{Hahn--Banach 定理}\label{ch:yosida-04}
\YosidaChapterPlaceholder{4}{The Hahn--Banach Theorems}{107--124}{102--118}

\section{实线性空间中的 Hahn--Banach 扩张定理}\label{sec:yosida-04-01}
\YosidaSectionPlaceholder{107}{102}
% TODO: 定义与记号；定理与完整证明；例子；习题；中英文核对。

\section{广义极限}\label{sec:yosida-04-02}
\YosidaSectionPlaceholder{109}{103}
% TODO: 定义与记号；定理与完整证明；例子；习题；中英文核对。

\section{局部凸的完备线性拓扑空间}\label{sec:yosida-04-03}
\YosidaSectionPlaceholder{110}{104}
% TODO: 定义与记号；定理与完整证明；例子；习题；中英文核对。

\section{复线性空间中的 Hahn--Banach 扩张定理}\label{sec:yosida-04-04}
\YosidaSectionPlaceholder{111}{105}
% TODO: 定义与记号；定理与完整证明；例子；习题；中英文核对。

\section{赋范线性空间中的 Hahn--Banach 扩张定理}\label{sec:yosida-04-05}
\YosidaSectionPlaceholder{112}{106}
% TODO: 定义与记号；定理与完整证明；例子；习题；中英文核对。

\section{非平凡连续线性泛函的存在性}\label{sec:yosida-04-06}
\YosidaSectionPlaceholder{113}{107}
% TODO: 定义与记号；定理与完整证明；例子；习题；中英文核对。

\section{线性映射的拓扑}\label{sec:yosida-04-07}
\YosidaSectionPlaceholder{115}{110}
% TODO: 定义与记号；定理与完整证明；例子；习题；中英文核对。

\section{嵌 \texorpdfstring{$X$}{X} 入其二次对偶 \texorpdfstring{$X''$}{X''}}\label{sec:yosida-04-08}
\YosidaSectionPlaceholder{117}{112}
% TODO: 定义与记号；定理与完整证明；例子；习题；中英文核对。

\section{几例对偶空间}\label{sec:yosida-04-09}
\YosidaSectionPlaceholder{119}{114}
% TODO: 定义与记号；定理与完整证明；例子；习题；中英文核对。

\section*{本章参考文献整理（占位）}
\phantomsection
\addcontentsline{toc}{section}{本章参考文献整理（占位）}
待核对原书本章参考文献。中译本印刷第 124 页起；后续按实际引用合并至本书参考文献表。
